\chapter{Análisis cualitativo (libro de códigos y matriz temática)}
\label{apx:cualitativo}

\section*{1. Criterio de análisis}
Las tres transcripciones canónicas se analizaron mediante un proceso de codificación temática. Los códigos se utilizaron para organizar recurrencias, tensiones y diferencias entre las perspectivas relevadas, conservando el contexto en el que se produjo cada observación. La codificación funciona como una herramienta de organización analítica y no como una cuantificación estadística de las entrevistas.

Para preservar la trazabilidad metodológica, se mantuvo la distinción entre una etapa de exploración (\textit{discovery}) y una instancia de retroalimentación posterior (\textit{feedback}) a la presentación de ImprontIA. En ENT-01 y ENT-02, la exploración inicial del problema se realizó antes de presentar la propuesta. En ENT-03, en cambio, se presentó una descripción general del proyecto al comienzo de la conversación, por lo que no se consideró que existiera un tramo de \textit{discovery} directamente comparable con los anteriores.

Convención de lectura: D = \textit{discovery}; F = \textit{feedback} posterior a la exposición de ImprontIA; ``Ambos'' indica presencia documentada en los dos momentos, sin trasladar una misma afirmación de un momento al otro.

El libro de códigos (\textit{codebook}) se reproduce tal como se utilizó en el análisis cualitativo. Los ejemplos incluidos corresponden a evidencia verificable en las transcripciones canónicas.

\section*{2. Libro de códigos}

\begin{center}\footnotesize
\begin{infobox}
\textbf{\textcolor{accent}{VAL-EVAL --- Integridad y validez de la evaluación}} \textit{(Momento: Ambos)}\\
\textit{Definición:} Preocupación por que el instrumento permita evaluar efectivamente aquello que se pretende evaluar y dentro de las condiciones establecidas.\\
\textit{Criterio de inclusión:} Autoría, cumplimiento de condiciones, correspondencia entre aprendizaje e instrumento.\\
\textit{Criterio de exclusión:} Problemas puramente técnicos sin efecto evaluativo.\\
\textit{Participantes:} ENT-01, ENT-02, ENT-03.\\
\textit{Evidencia:} ENT-01: ``La complejidad de evaluar a distancia sin una herramienta es enorme y, obviamente, atenta contra la integridad y la calidad académica.''
\end{infobox}
\vspace{2pt}
\begin{infobox}
\textbf{\textcolor{accent}{DIS-CONT --- Diseño y contexto de la evaluación}} \textit{(Momento: Ambos)}\\
\textit{Definición:} Necesidad de ajustar formato, consignas, recursos permitidos y grado de control al objetivo pedagógico y al contexto.\\
\textit{Criterio de inclusión:} Nivel de carrera, tipo de materia, pregunta, habilidades buscadas, uso permitido de IA.\\
\textit{Criterio de exclusión:} Supervisión aplicada de manera indiferenciada.\\
\textit{Participantes:} ENT-01, ENT-02, ENT-03.\\
\textit{Evidencia:} ENT-02: ``Entonces depende de la instancia de formación profesional, del enfoque de la evaluación, y ahí uno tiene que ajustar el uso de cada herramienta.''
\end{infobox}
\vspace{2pt}
\begin{infobox}
\textbf{\textcolor{accent}{SUP-MEC --- Mecanismos de supervisión}} \textit{(Momento: Ambos)}\\
\textit{Definición:} Métodos técnicos o procedimentales destinados a controlar las condiciones del examen.\\
\textit{Criterio de inclusión:} Cámara, micrófono, grabación de pantalla, bloqueo, cambio de pestaña, temporizadores.\\
\textit{Criterio de exclusión:} Juicio posterior acerca de culpabilidad.\\
\textit{Participantes:} ENT-01, ENT-02, ENT-03.\\
\textit{Evidencia:} ENT-02 describe ``plataformas que detectan si cambiás de pestaña'' y una evaluación con cámara y visualización de las manos.
\end{infobox}
\vspace{2pt}
\begin{infobox}
\textbf{\textcolor{accent}{FP-CONT --- Falsos positivos y contingencias}} \textit{(Momento: \textit{discovery})}\\
\textit{Definición:} Situaciones legítimas o fallas que pueden ser interpretadas por el sistema como incidentes.\\
\textit{Criterio de inclusión:} Ruido, movimientos naturales, notificaciones, fallos técnicos, terceros en el entorno.\\
\textit{Criterio de exclusión:} Conductas intencionalmente prohibidas comprobadas.\\
\textit{Participantes:} ENT-01, ENT-02.\\
\textit{Evidencia:} ENT-01: ``El docente revisa todos los resultados y valida cuáles son incidencias reales y cuáles son falsos positivos.''
\end{infobox}
\vspace{2pt}
\begin{infobox}
\textbf{\textcolor{accent}{FRIC-EST --- Estrés y fricción producidos por el control}} \textit{(Momento: \textit{discovery})}\\
\textit{Definición:} Presión adicional generada por las condiciones o mecanismos de supervisión.\\
\textit{Criterio de inclusión:} Temor a expulsión, necesidad de mantener postura/mirada, preocupación por alertas (\textit{flags}).\\
\textit{Criterio de exclusión:} Estrés propio del contenido del examen.\\
\textit{Participantes:} ENT-01, ENT-02.\\
\textit{Evidencia:} ENT-02: ``Por ahí no estás haciendo nada raro e igual te pone un flag de advertencia. Entonces te mete una presión que no solo va sobre el examen y los contenidos…''
\end{infobox}
\vspace{2pt}
\begin{infobox}
\textbf{\textcolor{accent}{PRIV-TRANS --- Privacidad y transparencia sobre la recolección}} \textit{(Momento: Ambos)}\\
\textit{Definición:} Preocupaciones por el acceso o registro de datos y necesidad de informar al evaluado qué se captura.\\
\textit{Criterio de inclusión:} Protección de datos, control del dispositivo, información previa.\\
\textit{Criterio de exclusión:} Reglas académicas no relacionadas con datos.\\
\textit{Participantes:} ENT-01, ENT-02.\\
\textit{Evidencia:} ENT-02: ``hay que avisarlo antes para que esa persona sepa qué datos se le están tomando''.
\end{infobox}
\vspace{2pt}
\begin{infobox}
\textbf{\textcolor{accent}{EVID-ACT --- Evidencia y accionabilidad}} \textit{(Momento: \textit{feedback} posterior)}\\
\textit{Definición:} Relación entre una señal detectada, la evidencia disponible y una acción académica defendible.\\
\textit{Criterio de inclusión:} Suficiencia probatoria, invalidación, repregunta, apelación, utilidad del reporte.\\
\textit{Criterio de exclusión:} Detección técnica considerada de manera aislada.\\
\textit{Participantes:} ENT-01, ENT-02, ENT-03.\\
\textit{Evidencia:} ENT-03: ``¿para qué voy a usar esa información y para qué puedo usarla? ¿Para qué voy a gastar recursos en obtener información que después no voy a poder utilizar?''
\end{infobox}
\vspace{2pt}
\begin{infobox}
\textbf{\textcolor{accent}{REV-HUM --- Revisión humana}} \textit{(Momento: Ambos)}\\
\textit{Definición:} Intervención de docentes u otros responsables para interpretar señales antes de concluir o actuar.\\
\textit{Criterio de inclusión:} Curaduría, interpretación contextual, validación de incidencias.\\
\textit{Criterio de exclusión:} Decisiones automáticas sin intervención humana.\\
\textit{Participantes:} ENT-01, ENT-02, ENT-03.\\
\textit{Evidencia:} ENT-01: ``Yo siento que la curaduría humana del resultado tiene que estar siempre''.
\end{infobox}
\vspace{2pt}
\begin{infobox}
\textbf{\textcolor{accent}{REGL-PREV --- Reglas y condiciones previas}} \textit{(Momento: \textit{feedback} posterior)}\\
\textit{Definición:} Definición y comunicación anticipada de conductas permitidas, consecuencias y procedimientos.\\
\textit{Criterio de inclusión:} Condiciones de aprobación, reglas de uso de recursos, consecuencias.\\
\textit{Criterio de exclusión:} Explicaciones dadas recién después de un incidente.\\
\textit{Participantes:} ENT-02, ENT-03.\\
\textit{Evidencia:} ENT-02: ``Tiene que estar de antemano, porque después aparece el `pero eso no lo decía en ningún lado'.''
\end{infobox}
\vspace{2pt}
\begin{infobox}
\textbf{\textcolor{accent}{VAL-COMP --- Validación complementaria / instancia oral}} \textit{(Momento: Ambos)}\\
\textit{Definición:} Uso de una segunda instancia para resolver dudas sobre conocimiento o autoría.\\
\textit{Criterio de inclusión:} Oral posterior, repregunta puntual, defensa de una respuesta, preguntas de validación.\\
\textit{Criterio de exclusión:} Oral como formato ordinario sin relación con una incertidumbre.\\
\textit{Participantes:} ENT-01, ENT-02, ENT-03.\\
\textit{Evidencia:} ENT-02: ``Si son pocos, para mí sería directamente una instancia oral… Si vos lo sabés y lo podés defender oralmente, pasa.''
\end{infobox}
\vspace{2pt}
\begin{infobox}
\textbf{\textcolor{accent}{CARGA-ESC --- Carga docente, adopción y escalabilidad}} \textit{(Momento: Ambos)}\\
\textit{Definición:} Trabajo adicional necesario para revisar reportes, aprender el sistema o actuar sobre grandes cantidades de alertas.\\
\textit{Criterio de inclusión:} Capacitación, autonomía, tiempo de revisión, cantidad de alumnos y de alertas.\\
\textit{Criterio de exclusión:} Complejidad técnica del algoritmo en sí misma.\\
\textit{Participantes:} ENT-01, ENT-02, ENT-03.\\
\textit{Evidencia:} ENT-01: ``Lo que queda pendiente ahora es formar al docente para que interprete correctamente los resultados.''
\end{infobox}
\vspace{2pt}
\begin{infobox}
\textbf{\textcolor{accent}{CTX-EST --- Condiciones concretas del estudiante}} \textit{(Momento: Ambos)}\\
\textit{Definición:} Restricciones materiales y ambientales que condicionan una evaluación remota.\\
\textit{Criterio de inclusión:} Dispositivos, notebooks laborales, familia, ruido, ambiente doméstico.\\
\textit{Criterio de exclusión:} Estrés puramente académico.\\
\textit{Participantes:} ENT-01, ENT-02, ENT-03.\\
\textit{Evidencia:} ENT-01: ``Otros tienen notebooks del trabajo y no las pueden usar porque tienen restricciones para instalar \textit{software}.''
\end{infobox}
\vspace{2pt}
\begin{infobox}
\textbf{\textcolor{accent}{AUT-DEC --- Límites a la automatización de decisiones}} \textit{(Momento: Ambos)}\\
\textit{Definición:} Frontera entre automatizar detección/flujo y delegar a una IA una decisión académica.\\
\textit{Criterio de inclusión:} Parametrización, intervención humana, disparadores automáticos controlados.\\
\textit{Criterio de exclusión:} Automatización meramente operativa que no afecta la evaluación.\\
\textit{Participantes:} ENT-01, ENT-03; parcialmente ENT-02.\\
\textit{Evidencia:} ENT-03: ``Que la inteligencia artificial no decida nada, por favor. Ahora, si yo lo puedo parametrizar, sí.''
\end{infobox}
\vspace{2pt}
\begin{infobox}
\textbf{\textcolor{accent}{CONS-CRIT --- Consistencia de criterios de actuación}} \textit{(Momento: \textit{feedback} posterior)}\\
\textit{Definición:} Necesidad de evitar respuestas completamente diferentes ante una misma señal según el docente.\\
\textit{Criterio de inclusión:} Protocolos, criterios homogéneos, relación indicador-acción.\\
\textit{Criterio de exclusión:} Diferencias pedagógicas legítimas entre materias.\\
\textit{Participantes:} ENT-02, ENT-03.\\
\textit{Evidencia:} ENT-03: ``Si no hay una indicación precisa de decir `si aparece este indicador, tenés que actuar de esta forma', quizás un profesor dice `no importa' y otro dice `esto es un error y hay que replantearlo'.''
\end{infobox}
\vspace{2pt}
\end{center}
Fuente: elaboración propia a partir de las tres transcripciones canónicas y del análisis cualitativo.

\section*{3. Matriz temática}
La matriz sintetiza la relación entre los temas cualitativos ya identificados y las perspectivas de los tres entrevistados. Se mantiene la distinción D/F y se evita convertir presencia temática en frecuencia estadística.

\begin{center}\footnotesize
\begin{infobox}
\textbf{\textcolor{accent}{Integridad / validez}} \textit{(D / F)}\\
\textbf{ENT-01:} D: considera que la evaluación remota requiere mecanismos que permitan controlar las condiciones.\\
\textbf{ENT-02:} D: la validez depende fuertemente de diseñar mejor la evaluación en contexto de IA.\\
\textbf{ENT-03:} F: vincula validez con la alineación entre objetivos, enseñanza, instrumento y recursos permitidos.\\
\textit{Síntesis analítica:} Los tres relacionan integridad con que el examen represente aquello que se pretendía evaluar. ENT-01 enfatiza control; ENT-02, diseño; ENT-03, alineación pedagógica y marco de evaluación.
\end{infobox}
\vspace{2pt}
\begin{infobox}
\textbf{\textcolor{accent}{Diseño y contexto}} \textit{(D / F)}\\
\textbf{ENT-01:} D: el control puede permitir evaluaciones digitales más cercanas a prácticas reales.\\
\textbf{ENT-02:} D: el nivel de la carrera, la consigna y el objetivo determinan si la IA debería permitirse y cuánto control tiene sentido.\\
\textbf{ENT-03:} F: calculadora, IA u otros recursos pueden ser válidos o inválidos dependiendo del objetivo del instrumento.\\
\textit{Síntesis analítica:} No aparece una defensa consistente del mismo nivel de control para todos los exámenes. ENT-02 explicita que en ciertos contextos avanzados un control estricto puede ser excesivo.
\end{infobox}
\vspace{2pt}
\begin{infobox}
\textbf{\textcolor{accent}{Mecanismos de supervisión}} \textit{(D / F)}\\
\textbf{ENT-01:} D: describe y valora una solución con cámaras y grabación de pantalla, aunque rechaza el modo de expulsión automática.\\
\textbf{ENT-02:} D: reconoce utilidad de bloqueo, cámaras, micrófono y temporizadores según el contexto.\\
\textbf{ENT-03:} F: considera que el monitoreo puede dar un marco más certero a las condiciones del examen.\\
\textit{Síntesis analítica:} Los tres admiten alguna utilidad posible del monitoreo; varía cuánto control consideran apropiado y con qué finalidad.
\end{infobox}
\vspace{2pt}
\begin{infobox}
\textbf{\textcolor{accent}{Falsos positivos / contingencias}} \textit{(discovery)}\\
\textbf{ENT-01:} D: los incidentes automáticos requieren distinguir casos reales de falsos positivos.\\
\textbf{ENT-02:} D: relata advertencias producidas por movimientos naturales de la mirada.\\
\textbf{ENT-03:} ---\\
\textit{Síntesis analítica:} ENT-01 y ENT-02 coinciden en que un evento técnico no equivale necesariamente a trampa. No corresponde afirmar consenso de los tres.
\end{infobox}
\vspace{2pt}
\begin{infobox}
\textbf{\textcolor{accent}{Estrés y fricción}} \textit{(D / F)}\\
\textbf{ENT-01:} D: la amenaza de expulsión agrega tensión y puede volver caótico el examen.\\
\textbf{ENT-02:} D/F: la supervisión puede obligar a conductas artificiales y generar presión; tras la presentación agrega problemas del entorno doméstico.\\
\textbf{ENT-03:} F: describe las dificultades del entorno doméstico para establecer un marco de examen.\\
\textit{Síntesis analítica:} Existe preocupación por efectos secundarios del control y del contexto remoto. ENT-03 enfatiza más el marco doméstico que el falso positivo algorítmico.
\end{infobox}
\vspace{2pt}
\begin{infobox}
\textbf{\textcolor{accent}{Privacidad / transparencia}} \textit{(D / F)}\\
\textbf{ENT-01:} D: identifica resistencia relacionada con protección de datos y registro del equipo.\\
\textbf{ENT-02:} F: considera necesario avisar previamente qué datos se recopilan.\\
\textbf{ENT-03:} ---\\
\textit{Síntesis analítica:} ENT-01 y ENT-02 señalan la dimensión de datos del estudiante. No existe evidencia directa suficientemente clara para afirmar consenso de los tres.
\end{infobox}
\vspace{2pt}
\begin{infobox}
\textbf{\textcolor{accent}{Evidencia y accionabilidad}} \textit{(feedback)}\\
\textbf{ENT-01:} F: cuestiona sancionar a partir de eventos sin evidencia que permita sostener la acusación.\\
\textbf{ENT-02:} F: propone usar una alerta para seleccionar a quién repreguntar o evaluar oralmente.\\
\textbf{ENT-03:} F: cuestiona el valor de obtener información si luego no habilita ninguna acción clara.\\
\textit{Síntesis analítica:} Después de conocer la propuesta, los tres desplazan la pregunta desde detectar hacia qué se hace con lo detectado. ENT-01 exige evidencia más contundente para sanciones; ENT-02 y ENT-03 exploran usos complementarios.
\end{infobox}
\vspace{2pt}
\begin{infobox}
\textbf{\textcolor{accent}{Revisión humana / automatización}} \textit{(D / F)}\\
\textbf{ENT-01:} D: afirma que la curaduría humana debe permanecer.\\
\textbf{ENT-02:} D/F: usa criterio docente caso a caso y propone validación humana posterior.\\
\textbf{ENT-03:} F: rechaza que una IA decida autónomamente, aunque acepta comportamientos previamente parametrizados.\\
\textit{Síntesis analítica:} La decisión académica no aparece como algo que deba delegarse ciegamente a una detección automática. ENT-03 distingue autonomía de IA de automatización parametrizada.
\end{infobox}
\vspace{2pt}
\begin{infobox}
\textbf{\textcolor{accent}{Instancia oral / validación}} \textit{(D / F)}\\
\textbf{ENT-01:} D: informa que anteriormente se validaban evaluaciones mediante examen oral, aunque lo considera insuficiente como único control.\\
\textbf{ENT-02:} F: la propone como respuesta dirigida a casos señalados.\\
\textbf{ENT-03:} F: propone oral, repreguntas o preguntas de validación como posibles complementos.\\
\textit{Síntesis analítica:} La instancia complementaria aparece como mecanismo para obtener evidencia adicional. ENT-01 cuestiona su suficiencia como solución primaria; ENT-02 y ENT-03 la ven útil como segunda instancia.
\end{infobox}
\vspace{2pt}
\begin{infobox}
\textbf{\textcolor{accent}{Reglas previas / consistencia}} \textit{(feedback)}\\
\textbf{ENT-01:} ---\\
\textbf{ENT-02:} F: insiste en explicitar reglas y consecuencias antes del examen.\\
\textbf{ENT-03:} F: advierte que docentes distintos podrían actuar de manera diferente frente al mismo indicador.\\
\textit{Síntesis analítica:} ENT-02 y ENT-03 destacan reglas procedimentales. ENT-02 pone el foco en información previa al alumno; ENT-03, en consistencia entre docentes.
\end{infobox}
\vspace{2pt}
\begin{infobox}
\textbf{\textcolor{accent}{Carga y escalabilidad}} \textit{(D / F)}\\
\textbf{ENT-01:} D: identifica capacitación e interpretación de reportes como pendiente operativo.\\
\textbf{ENT-02:} F: plantea el problema de tener, por ejemplo, 80 alertas en un curso de 100 alumnos.\\
\textbf{ENT-03:} F: cuestiona invertir recursos en información que no pueda transformarse en acción.\\
\textit{Síntesis analítica:} Una herramienta genera trabajo posterior y ese trabajo condiciona su utilidad: capacitación, volumen de casos y costo/beneficio de la información.
\end{infobox}
\vspace{2pt}
\end{center}
Fuente: elaboración propia a partir de la matriz temática del análisis cualitativo.

\section*{4. Nota metodológica de lectura}
La presencia de un código o tema en más de una entrevista no se interpreta como frecuencia estadística ni como medida de consenso. Las entrevistas aportan perspectivas individuales situadas; cuando una idea aparece solo en dos participantes, se conserva ese alcance y no se atribuye al conjunto de los tres.

ENT-01 y ENT-02 cuentan con un tramo de \textit{discovery} previo a la presentación de ImprontIA. ENT-03 no cuenta con un \textit{discovery} limpio directamente comparable. En consecuencia, los aportes producidos en \textit{feedback} se mantienen diferenciados y no se presentan como necesidades espontáneamente descubiertas.

Las opiniones de los entrevistados se atribuyen a cada participante y no a UADE como institución. El análisis cualitativo documenta experiencias, tensiones y criterios relevados; no constituye una validación directa de ImprontIA.
